Section~\ref{sec:lim} names two assumptions and says neither is tested: the
moving-block construction needs stationarity and mixing that an event log is
not guaranteed to have, and the same assumption is what makes
Definition~\ref{def:estimand}'s population limit well defined under a temporal
split. This section tests the first of them, on every pair, and reports the
answer whichever way it comes out.

\paragraph{The diagnostic, and the null it is read against}
At the reference cell the model is fitted \emph{once} on the training half and
then evaluated on the test half cut into \nDriftBlocks\ consecutive blocks in
time. Under a stationary process the per-block increment should differ only by
sampling error. A block is a fifth of the test half, so it carries more
sampling error than the whole and \emph{some} spread is guaranteed; comparing
the spread with the pair's own interval, which is a whole-sample object, would
therefore call ordinary noise drift. The reference is instead a
\textbf{stationarity null}: the same test rows re-partitioned at random into
blocks of the same sizes, \nDriftPerm\ times, which destroys the time order
and nothing else. The observed spread is reported as a percentile of that
reference.

\begin{table}[t]
\centering
\small
\resizebox{\ifdim\width>\linewidth\linewidth\else\width\fi}{!}{%%
\begin{tabular}{llllllllll}
\toprule
log & target & V & range over blocks & spread & null median & percentile & trend & sign varies & drifts \\
\midrule
BPIC13\textbackslash \_closed & handover & $-$0.009 & [$-$0.044, +0.011] & 0.055 & 0.054 & 51.5 & 0.50 & yes & no \\
BPIC13\textbackslash \_incidents & duration & 0.028 & [$-$0.003, +0.055] & 0.058 & 0.043 & 76.8 & $-$0.50 & yes & no \\
BPIC13\textbackslash \_incidents & handover & 0.037 & [+0.032, +0.042] & 0.010 & 0.022 & 3.3 & $-$0.50 & no & no \\
BPIC14 & duration & 0.096 & [+0.069, +0.112] & 0.043 & 0.021 & 98.0 & 0.10 & no & yes \\
BPIC14 & handover & 0.167 & [+0.157, +0.184] & 0.027 & 0.047 & 13.5 & 0.00 & no & no \\
BPIC15\textbackslash \_1 & duration & 0.001 & [$-$0.077, +0.103] & 0.179 & 0.069 & 100.0 & $-$0.10 & yes & yes \\
BPIC15\textbackslash \_1 & handover & 0.007 & [+0.002, +0.012] & 0.010 & 0.021 & 6.8 & 0.10 & no & no \\
BPIC15\textbackslash \_3 & duration & 0.035 & [$-$0.013, +0.038] & 0.051 & 0.068 & 22.5 & 0.10 & yes & no \\
BPIC15\textbackslash \_3 & handover & $-$0.003 & [$-$0.053, +0.049] & 0.102 & 0.074 & 83.5 & 0.00 & yes & no \\
BPIC15\textbackslash \_4 & duration & 0.035 & [$-$0.036, +0.043] & 0.079 & 0.048 & 92.8 & $-$0.40 & yes & no \\
BPIC15\textbackslash \_4 & handover & 0.001 & [$-$0.041, +0.020] & 0.062 & 0.038 & 92.0 & 0.67 & yes & no \\
BPIC15\textbackslash \_5 & duration & 0.032 & [+0.018, +0.107] & 0.089 & 0.106 & 35.5 & 0.60 & no & no \\
BPIC17 & duration & 0.002 & [$-$0.010, +0.011] & 0.022 & 0.015 & 86.5 & $-$0.60 & yes & no \\
BPIC19 & duration & 0.024 & [$-$0.037, +0.063] & 0.100 & 0.010 & 100.0 & $-$0.80 & yes & yes \\
Helpdesk & duration & $-$0.020 & [$-$0.103, +0.015] & 0.118 & 0.042 & 100.0 & 0.30 & yes & yes \\
RoadFines & duration & 0.031 & [+0.025, +0.069] & 0.044 & 0.012 & 100.0 & 1.00 & no & yes \\
Sepsis & duration & 0.153 & [+0.106, +0.230] & 0.124 & 0.105 & 64.5 & 0.60 & no & no \\
UCI498 & duration & $-$0.007 & [$-$0.015, +0.009] & 0.025 & 0.013 & 99.0 & 0.10 & yes & yes \\
UCI498 & handover & $-$0.003 & [$-$0.005, $-$0.002] & 0.003 & 0.005 & 9.5 & $-$0.60 & no & no \\
\bottomrule
\end{tabular}%
}
\caption{IS THE INCREMENT STATIONARY ACROSS THE TEST HALF? The reference cell's model is fitted once on the training half and evaluated on the test half cut into \nDriftBlocks\ consecutive blocks in time. `spread' is the range of the per-block increment; `null median' is the median spread over \nDriftPerm\ RANDOM re-partitions of the same rows into blocks of the same sizes, which destroys the time order and nothing else; `percentile' places the observed spread in that reference. A block is a fifth of the test half and so carries more sampling error than the whole, which is why the comparison is against the null and not against the pair's own interval. `trend' is the Spearman correlation of the increment with the block index: it separates drift, which has an order, from heterogeneity, which does not. `drifts' marks a percentile in the null's upper tail at the declared level $\alpha = 0.05$.}
\label{tab:drift}
\end{table}


\paragraph{What it shows}
\textbf{On \nDrifts\ of \nDriftPairs\ pairs the increment is not stationary
across the test half} --- the time-ordered spread sits in the upper tail of
the null at the declared level $\alpha = 0.05$ --- and on
\nDriftAboveNullMedian\ it sits above the null's median. The median observed spread is \driftSpreadMedian\ AUC,
\driftOverNullMedian\ times the null's median, and on the worst pair it is
\driftOverNullMax\ times. On \nDriftSignVaries\ of \nDriftPairs\ pairs the
\emph{sign} of the increment varies between blocks of the same test half.

\textbf{It is heterogeneity rather than trend.} Only \nDriftTrend\ of
\nDriftPairs\ pairs shows a Spearman correlation between the increment and the
block index that reaches conventional significance, so what the diagnostic
finds is not a register decaying smoothly with time: it is a test half whose
periods are not exchangeable. Two candidate mechanisms are in the table.
Prevalence moves by a median \driftPrevSpreadMedian\ between the first and
last block, and by \driftPrevSpreadMax\ on the worst pair; and the share of
test rows whose register value was never seen in training rises from a median
\driftUnseenFirstPct\ in the first block to \driftUnseenLastPct\ in the last.

\paragraph{What follows}
Two things, and neither is a correction. First, a pair that fails this test is
not a pair whose number is wrong: it is a pair whose estimand is a moving
target, which is a fact a reader should be given rather than an error to be
repaired --- and it is the paper's own thesis applied to time, which
$\mathcal{S}$ does not contain as an axis. Second, the interval and the band
do not cover it. Both are computed under an assumption this diagnostic finds
violated on \nDrifts\ pairs, so on those pairs the interval understates what a
reader would see on a different stretch of the same log, and the direction of
the understatement is not signed. Adding time as a declared axis --- an
evaluation period, with the increment reported per period --- is the change
this measurement implies, and it is not made here.
